% from ega4-9-fr.tex
\oldpage[IV-3]{89}

contenant $x$, on a
\[
\operatorname{codim}_x(Z_s,X_s)=
\inf_\beta\left(\sup_\alpha
\bigl(\operatorname{codim}(Z_s^{(\beta)},X_s^{(\alpha)})\bigr)\right),
\]
$(\alpha,\beta)$ variant dans l'ensemble des couples tels que
$x\in Z_s^{(\beta)}\subset X_s^{(\alpha)}$ ($\mathbf 0$, 14.2.6). Comme les
préschémas algébriques irréductibles sont biéquidimensionnels (5.2.1), on peut
écrire, en vertu de ($\mathbf 0$, 14.3.3.1)
\begin{equation}
\tag{9.9.1.1}
\label{IV.9.9.1.1-fr}
\operatorname{codim}_x(Z_s,X_s)=
\inf_\beta\left(\sup_\alpha
\bigl(\dim(X_s^{(\alpha)})-\dim(Z_s^{(\beta)})\bigr)\right)
\end{equation}
avec le même choix des couples $(\alpha,\beta)$. Comme $h_s$ est fidèlement
plat et quasi-compact, pour tout couple formé d'une composante irréductible
$(X_0)_{s_0}^{(\lambda)}$ de $(X_0)_{s_0}$ contenant $x_0$ et d'une composante
irréductible $(Z_0)_{s_0}^{(\mu)}$ de $(Z_0)_{s_0}$ contenant $x_0$ et contenue
dans $(X_0)_{s_0}^{(\lambda)}$, il existe un couple
$(X_s^{(\alpha)},Z_s^{(\beta)})$ du type décrit plus haut et tel que
$Z_s^{(\beta)}$ domine $(Z_0)_{s_0}^{(\mu)}$ et que $X_s^{(\alpha)}$ domine
$(X_0)_{s_0}^{(\lambda)}$ (2.3.5). La formule (9.9.1.1) (et la formule analogue
appliquée à $(X_0)_{s_0}$) montrent alors, en vertu de (4.2.7), que l'on a
\[
\operatorname{codim}_x(Z_s,X_s)=
\operatorname{codim}_{x_0}((Z_0)_{s_0},(X_0)_{s_0}).
\]

On voit ainsi que si $E$ (resp. $E_0$) est l'ensemble des $x\in X$ (resp. des
$x_0\in X_0$) vérifiant l'une des conditions (i), (ii), (iii) de l'énoncé
(resp. la même condition), on a $E=h^{-1}(E_0)$, et en vertu de (1.8.2), on
voit qu'on peut se borner au cas où $A$ est \emph{noethérien}, donc aussi $B$.
Compte tenu de ($\mathbf 0_{\mathrm{III}},9.2.3$), ainsi que de (9.9.1.1), on
est ramené à voir que pour tout $x\in X$, il y a un voisinage $V$ de $x$ dans
$\overline{\{x\}}$ tel que, pour tout $x'\in V$, l'ensemble des dimensions des
composantes irréductibles de $X_{f(x')}$ (resp. $Z_{f(x')}$) contenant $x'$ est
\emph{le même}, et en outre qu'il en est de même de l'ensemble des couples
\[
\bigl(\dim(Z_{f(x')}^{(\beta)}),\dim(X_{f(x')}^{(\alpha)})\bigr)
\]
pour les couples formés d'une composante irréductible $X_{f(x')}^{(\alpha)}$
de $X_{f(x')}$ et d'une composante irréductible $Z_{f(x')}^{(\beta)}$ de
$Z_{f(x')}$ contenue dans $X_{f(x')}^{(\alpha)}$ et contenant $x'$. On peut
évidemment pour cela remplacer $S$ par le sous-préschéma réduit $S'$ de $S$
ayant $\overline{\{f(x)\}}$ pour espace sous-jacent, et $X$ par
$X'=f^{-1}(S')$, les fibres de $X$ et $X'$ aux points de $S'$ étant les mêmes.
Autrement dit, on peut se borner au cas où $S$ est \emph{intègre} et où
$\eta=f(x)$ est son point générique.

Par hypothèse $Z_\eta$ est fermé dans $X_\eta$~; comme $Z$ est constructible,
il résulte de (8.3.11), appliqué suivant la méthode de (8.1.2, a)), que l'on
peut, en remplaçant au besoin $S$ par un voisinage ouvert de $\eta$, supposer
que $Z$ est \emph{fermé} dans $X$. Soient $X_i$ (resp. $Z_j$) les composantes
irréductibles de $X$ (resp. $Z$) contenant $x$~; en vertu de
($\mathbf 0_{\mathrm I},2.1.8$), les $X_i\cap X_\eta$ (resp.
$Z_j\cap X_\eta$) sont les composantes irréductibles de $X_\eta$ (resp.
$Z_\eta$) contenant $x$~; en vertu de (9.5.1), on peut supposer en outre, en
restreignant au besoin $S$ à un voisinage de $\eta$, que les $X_i$ (resp.
$Z_j$) sont exactement les composantes irréductibles de $X$ (resp. $Z$)
rencontrant $\overline{\{x\}}$ et que les
$(X_i)_s\cap\overline{\{x\}}$ et $(Z_j)_s\cap\overline{\{x\}}$ sont non vides
pour tout $s\in S$. Cela étant, il résulte encore de (9.7.1) que l'on peut
supposer, en restreignant $S$, que les couples $(i,j)$ tels que
$(Z_j)_s\subset(X_i)_s$ sont les mêmes pour tout $s\in S$. La conclusion
résulte alors de (9.5.6)~: pour tout $s$ assez voisin de $\eta$, toutes les
composantes irréductibles de $(X_i)_s$ (resp. $(Z_j)_s$) ont une même dimension
égale à celle de $(X_i)_\eta$ (resp. $(Z_j)_\eta$). En outre, si $(i,j)$ est un
couple tel que $Z_j\not\subset X_i$, $(X_i)_\eta$ ne contient pas le

